\chapter*{Introduction: The Question Under the Question}
\addcontentsline{toc}{chapter}{Introduction: The Question Under the Question}
\markboth{INTRODUCTION}{INTRODUCTION}

The easiest answer to \emph{How did you get rich?} is a number. It is also the
answer most likely to confuse us.

A person says that a company sold for a billion dollars. We imagine a billion
dollars arriving in one bank account. The transaction may instead include
stock, debt, earn-outs, other owners, dilution, taxes, and a price paid for the
whole company rather than for one person's share. Another person reports a
hundred million dollars in annual revenue. Revenue sounds like wealth until
wages, inventory, rent, interest, refunds, taxes, and the cash needed for next
month enter the room. A third person owns shares worth tens of millions at
today's market price. That is wealth of a kind, but it is not the same as cash
that can be spent without changing the price, control, or tax bill.

The number attracts us because it seems to settle the question. It lets the
interviewer and the audience rank lives quickly. Yet every serious answer in
these conversations begins when the ranking breaks down.

One aviation entrepreneur defines being rich relative to what a life costs.
Grant Cardone takes a million-dollar balance, spreads it across the remaining
years of a hypothetical life, and uses the result to provoke the listener: a
large-looking pile may support much less freedom than its owner imagines. Glenn
Boyd, looking back on childhood work in the Texas heat and alongside his mother
washing dishes, says money mattered greatly when it removed deprivation. He
also says that beyond a certain point, more luxurious versions of the same
consumption stopped changing the experience. Tony Stephens tears a strip of
paper into pieces to show that the portion already lived cannot be repurchased.

These are not four versions of one formula. They are four attacks on the same
mistake: treating the stock of money as if it were the whole of wealth.

\section*{The categories hidden inside ``rich''}

We need a small vocabulary before the stories can teach us anything. Earned
income pays for current labor. Business revenue records sales before most
obligations. Profit is an accounting result under a defined set of rules. Cash
flow records when money actually moves. Equity is a residual claim on what an
asset or company may be worth after other claims. Liquidity describes how
readily value can become spendable cash without a damaging concession. Security
is the ability to absorb plausible shocks. Freedom is the ability to choose
among meaningful alternatives.

These categories touch one another, but none can stand in for all the rest. A
high income paired with high fixed spending can create a fragile life. An owner
can be wealthy on paper and unable to meet payroll. A modest household can have
little status and considerable freedom because its obligations are low, its
skills are portable, and its relationships are strong. A founder can control a
valuable company and still be unable to leave it for a week.

The distinction matters because the book follows a sequence. Labor can create
income. A margin between income and spending creates choice. Useful work can
create customer value. Repeated customer value can support a company. Rights,
equity, and systems can allow the creator to participate beyond the next hour
worked. Capital can preserve and enlarge those claims. None of this determines
whether the resulting life is good. That judgment has to govern the sequence
from the beginning.

\section*{Experience is evidence, not a law}

Successful people often explain success through the traits that remain visible
to them: persistence, courage, focus, faith, speed, salesmanship, or attention
to detail. Those traits matter. The interviews contain repeated examples of a
person making another call, learning a technical field, staying with an
unfashionable business, or acting while an opportunity remained open.

The same stories contain less flattering causes. Timing matters. One person
bought before a market rose and readily calls the result luck. Another held a
concentrated position while its value fell by millions. A founder benefited
from a partner's capital at the moment his own money ran out. A professional
credential opened one path while a chance encounter opened another. Family
labor, migration, access to a wealthy place, a healthy body, a tolerant spouse,
an institution willing to lend, and a customer prepared to trust a newcomer
all alter what persistence can accomplish.

We do not honor a person's agency by erasing those conditions. We make the
lesson more useful by asking a better question: \emph{What mechanism operated
here, under what conditions, and what would cause it to fail?}

That question changes the tone of the book. ``Never quit'' becomes a discussion
of encounter probability, evidence, health, and stopping rules. ``Use debt''
becomes a cash-flow stress test. ``Build a personal brand'' becomes a path from
attention through delivery to remembered trust. ``Find a great partner''
becomes a set of reversible tests, decision rights, capital obligations, and
exit terms. The vivid sentence survives, but it has to carry its own weight.

\section*{The argument ahead}

The book has five movements. The first asks what money is for and what a reader
is actually starting with. The second follows value outward: problems,
customers, products, price, sales, and distribution. The third follows value
into ownership: equity, systems, recurring revenue, teams, rights, control, and
succession. The fourth slows the story down. It separates the numbers, examines
debt and property, studies compounding and concentration, and builds the
reserves and protective structures that keep one error from becoming ruin. The
fifth returns to the human being who wanted freedom in the first place.

The governing idea is simple enough to remember and demanding enough to test
every chapter:

\begin{equation}
\begin{aligned}
\text{a well-lived life}&=\text{the end},\\
\text{wealth}&=\text{infrastructure},\\
\text{leadership}&=\text{stewardship}.
\end{aligned}
\end{equation}

Wealth is infrastructure because it can support time, safety, care, mobility,
experimentation, recovery, and refusal. Infrastructure is not sacred. It is
judged by what it allows and by whom it burdens. Leadership is stewardship
because no durable company is made by the visible owner alone. Customers trust
it, workers carry it, families absorb its uncertainty, lenders and communities
permit it to operate, and future people inherit what it builds or damages.

The title of this book is deliberately addressed to someone else: \emph{How
you got rich.} Curiosity begins there. The useful work begins when the pronoun
changes. What are \emph{we} trying to make possible? What must remain true if
the plan succeeds? What are we unwilling to spend, even for more?

Those questions cannot be postponed until the final chapter. They are the
first form of financial discipline.
